This study develops a unified framework for the three-dimensional large-deformation dynamics of a magnetically actuated flexible catheter subject to magnetic, fluid, and vessel-wall contact effects in microgravity. A geometrically exact Cosserat-rod model, a joint translational--rotational Galerkin discretization, local constraint reduction, and a consistent Legendre transformation provide a finite-dimensional canonical Hamiltonian system and COMSOL-compatible state and port interfaces. The fixed \(N_{\mathrm{diag}}=2\) diagnostic basis uses the paired-coordinate contract \(q=\Pi_{\mathrm{conf}}(U)\), \(p=B_N(q)^TP_{\mathrm{full}}\), \(\dot q=M_N(q)^{-1}p\), and \(\dot p=\dot B_N(q)^TP_{\mathrm{full}}+B_N(q)^T\dot P_{\mathrm{full}}\). It is used to test consistency of the data mapping and does not establish \(N_{\mathrm{diag}}=2\) as the final sufficient modal order. All three candidate trajectories in the COMSOL--Galerkin bridge satisfy the independent \(2\%\) force-v2 joint-L2 criterion, whereas the original \(0.5\%\) strict perturbation-energy criterion is not satisfied overall; the two criteria are not interchangeable. The low-dimensional P0 study and the three-trajectory P1/P1.5/P1.6 studies show a conditional advantage in the derivative-supervised pilot, substantial gray-box residual gains over physics-only without stable one-step-field superiority, and a marked loss of predictive accuracy when exact equilibrium constraints are imposed, respectively. After correcting a cross-trajectory comparison, isolating the error source, identifying the cotangent projection contract, comparing solver-output semantics, and correcting fixed-basis initialization, a three-level time-step comparison yields observed orders of \(1.64186\)--\(2.00088\). For the \(1\,\mathrm{s}\) amplitude-sentinel case, G2, pointwise G3, impulse G3, and peak-to-peak perturbation-energy drift are \(0.00146226\), \(0.000221162\), \(0.0000563606\), and \(0.00307057<0.005\), respectively. The finite evidence supports at most this single diagnostic sentinel as a candidate input to a subsequent A1 qualification experiment; it does not certify data-generation qualification across modes, amplitudes, or operating conditions, and it neither establishes BDF symplecticity nor constitutes numerical certification of P2 or KAM conclusions. The planned 96-trajectory data set, full-scale P2 training, open-port and high-fidelity studies, numerical invariant tori, and engineering safe domain remain subjects of future work.
